\documentclass{article} % say
\usepackage{amsmath}
\usepackage{gensymb}
\usepackage{hyperref}
\usepackage{tikz}
\usepackage{tkz-euclide}
\usepackage{amsfonts}
\usetikzlibrary {babel,angles,quotes,positioning}
\usepackage[utf8]{inputenc}
\usepackage[russian]{babel}
\begin{document}

\textbf{Задача}: В соревновании участвуют $25$ команд, пронумерованных от $1$ до $25$. Соревнование представляет собой серию раундов, в каждом из которых встречаются две команды. Команда с б\'{о}льшим номером всегда выигрывает раунд. В первом раунде две команды выбираются равновероятно из всех возможных команд. После каждого раунда проигравшая команда выбывает из соревнования, а выигрывшая выходит в следующий раунд, где встречается со случайной командой, выбранной равновероятно из ещё оставшихся в соревновании. Известно, что команда, выбранная в первом раунде выиграла $18$ раундов подряд (включая самый первый). С какой вероятностью она выиграет $19$-й раунд?

\clearpage

\textbf{Длинное решение}: Обозначим за $A$ событие, что команда выиграла первые $19$ раундов, а за $B$~--- выиграла первые $18$. Тогда мы хотим посчитать $P(A|B)$. В нашем случае $P(A|B) = \frac{P(A)}{P(B)}$, так как $A \subset B$. Если команда выиграла $k$ раундов, то её номер как минимум $k$ (должна найтись $k-1$ команда с меньшим номером).

\begin{equation}
  P(A) = \sum\limits_{i=19}^{25} {\color{blue} \frac{2}{25} } \prod\limits_{j=1}^{19} \dfrac{\color{red} i - j}{\color{brown} 25 - j} = \frac{1}{10}
\end{equation}

\begin{equation}
  P(B) = \sum\limits_{i=18}^{25} {\color{blue} \frac{2}{25} } \prod\limits_{j=1}^{18} \dfrac{\color{red} i - j}{\color{brown} 25 - j} = \frac{2}{19}
\end{equation}

Синим цветом в формулах обозначается вероятность, что команда с номером $i$ была выбрана в первый раунд. Затем $j$ можно воспринимать как номер раунда. Числитель (выделенный красным) означает количество команд с номером меньшим $i$, которые ещё не выбыли из соревнования перед раундом $j$. Знаменатель (выделенный коричневым) обозначает количество команд, которые ещё не выбыли из соревнования перед раундом $j$ исключая команду $i$.

\textbf{Ответ}: $\frac{19}{20}$

\clearpage

\textbf{Короткое решение}: Обозначим за $f_n^m$ вероятность того, что среди $n$ команд, команда, выигравшая первый раунд, выиграла ещё как минимум $m-1$. Нас спрашивают $\frac{f_{25}^{19}}{f_{25}^{18}}$.

Поставим в соответствие соревнованию перестановку из $n$ чисел: первые два числа соответсвуют командам, игравшим в первом раунде, третье~--- команде, которая встретилась с победителем первого раунда, и так далее. Легко видеть, что команда из первого раунда выигрывает первые $m$ раундов, в случае если максимум первых $m+1$ элементов перестановки находится на первой или второй позиции. Таким образом, $f_{n}^m = \frac{2}{m+1}$.

\textbf{Ответ}: $\frac{f_{25}^{19}}{f_{25}^{18}} = \frac{19}{20}$

\textbf{Источник}: \url{https://t.me/obznam/1423}

\end{document}


